\section{Further directions}\label{sec:further}

The restriction map connects three parts of the transition: its kernel is
the surviving second cohomology, its image gives the integral relations
among cohomological vanishing cycles, and its periods give the logarithmic
spectral energy. The polarized harmonic quotient supplies the mass.
The split-quintic calculation makes the associated integral rings and
attachments explicit; the toric examples show how a higher-dimensional
quotient produces several distinct coefficients.

For an actual base change $s=t^k$ of a family satisfying the first-order
hypotheses, the polarized fibre at $t$ is exactly the old fibre at $s=t^k$.
Ricci-flat uniqueness therefore gives
\[
 \lambda^+_{2,j}(t)=\frac{2\pi\mu_j}{k|\log t|}(1+o(1)).
\]
Counts, gaps and type are inherited, as is the rank-one one-sided remainder
for fixed $k$. Equal local smoothing orders alone do not imply that an
arbitrary global family descends through this base change. Different node
orders or ratios approaching zero require metric estimates uniform on the
enlarged parameter domain; the separate cap capacities do not supply them.

The proof gives a leading coefficient for each fixed family. A two-sided
next-order expansion would require quantitative control of the limiting
spectral subspaces and the full harmonic projection. Uniformity as the
quotient spectral matrix itself degenerates is a further problem.
The small two-form modes have exact three-form partners under $d$, and
Hodge duality gives related modes in degree four; the argument does not
exhaust the middle-degree small spectrum.

The tangent results in Appendix~\ref{sec:tangent} concern connections
induced by the Ricci-flat geometry. Extending the analysis to a separately
specified bundle would require compatible deformation and stability data,
uniform bundle inverses, and control of possible curvature concentration.
Such bundle spectra and HYM flows are separate from the Hodge theorem
proved here.
